% GENERATED FILE. Edit canonical lessons, then run npm run generate:books.
% canonical-source-hash: fnv1a64:acada4e5e46fb2b8
% canonical-lessons: JA-C416-yoake, JA-C416-hinode, JA-C416-yuuyake, JA-C416-hiekomu, JA-C416-terasu

\chapter{Dawn and Dusk}
\label{ch:ja-dawn-and-dusk}
\hlchaptermodality{416}

\begin{chapteropening}
\textbf{By the end of this chapter:} \emph{I can talk about sunrise, sunset and the cold of the night.}

Complete the last lesson of chapter 416: I can talk about sunrise, sunset and the cold of the night.
\end{chapteropening}

\section[\texorpdfstring{\ja{よあけ} (yoake)}{よあけ (yoake)}]{\ja{よあけ} (yoake) — dawn; daybreak}

\label{lesson:JA-C416-yoake}



\begin{quote}
Before the new one: say the Japanese for to get something wet, then the Japanese for sheltering from the rain.
\end{quote}

\subsection*{You'll want to know: \ja{よあけ}}
\textbf{\ja{よあけ}} — \emph{yoake} — \textquotedblleft{}dawn; daybreak\textquotedblright{}.

\begin{quote}
\textbf{\ja{よあけに} \ja{やどを} \ja{でました。}} — \emph{yoake ni yado o demashita} — \textquotedblleft{}We left the inn at dawn.\textquotedblright{}
\end{quote}

\begin{cousinweb}
\textbf{\ja{よあけ}} (\emph{yoake}) is a native Japanese word (wago), inherited rather than borrowed from Chinese, built from \emph{yo}, \textquotedblleft{}night\textquotedblright{} and \emph{ake}, \textquotedblleft{}opening\textquotedblright{} (from \emph{akeru}); its kanji were attached later as kun'yomi (JMdict).
\end{cousinweb}

\begin{grammarlens}[title={What you've built}]
One more word for getting around.
\end{grammarlens}

\subsection*{Your turn}
\begin{itemize}
\raggedright
  \item \emph{Say it:} \emph{yoake}
  \item \emph{Say it:} \emph{yoake}, once more
  \item \emph{Say it:} \emph{yoake ni yado o demashita} — We left the inn at dawn.
  \item \emph{Recall:} say the Japanese for to get something wet, then the Japanese for sheltering from the rain, then say \emph{yoake} again
\end{itemize}

\begin{tcolorbox}[breakable,colback=teal!4,colframe=teal!35!black,title={Before you move on}]
What does \ja{よあけ} mean? (\textquotedblleft{}Dawn; daybreak\textquotedblright{}.) Say it once more.
\end{tcolorbox}
\section[\texorpdfstring{\ja{ひので} (hinode)}{ひので (hinode)}]{\ja{ひので} (hinode) — sunrise}

\label{lesson:JA-C416-hinode}



\begin{quote}
Before the new one: say the Japanese for sheltering from the rain, then the Japanese for dawn.
\end{quote}

\subsection*{You'll want to know: \ja{ひので}}
\textbf{\ja{ひので}} — \emph{hinode} — \textquotedblleft{}sunrise\textquotedblright{}.

\begin{quote}
\textbf{\ja{やまの} \ja{うえで} \ja{ひのでを} \ja{まちました。}} — \emph{yama no ue de hinode o machimashita} — \textquotedblleft{}We waited for the sunrise on the mountain.\textquotedblright{}
\end{quote}

\begin{cousinweb}
\textbf{\ja{ひので}} (\emph{hinode}) is a native Japanese word (wago), inherited rather than borrowed from Chinese, built from \emph{hi}, \textquotedblleft{}sun\textquotedblright{}, the particle \emph{no} and \emph{de}, \textquotedblleft{}coming out\textquotedblright{} (from \emph{deru}); its kanji were attached later as kun'yomi (JMdict).
\end{cousinweb}

\begin{grammarlens}[title={What you've built}]
One more word for getting around.
\end{grammarlens}

\subsection*{Your turn}
\begin{itemize}
\raggedright
  \item \emph{Say it:} \emph{hinode}
  \item \emph{Say it:} \emph{hinode}, once more
  \item \emph{Say it:} \emph{yama no ue de hinode o machimashita} — We waited for the sunrise on the mountain.
  \item \emph{Recall:} say the Japanese for sheltering from the rain, then the Japanese for dawn, then say \emph{hinode} again
\end{itemize}

\begin{tcolorbox}[breakable,colback=teal!4,colframe=teal!35!black,title={Before you move on}]
What does \ja{ひので} mean? (\textquotedblleft{}Sunrise\textquotedblright{}.) Say it once more.
\end{tcolorbox}
\section[\texorpdfstring{\ja{ゆうやけ} (yūyake)}{ゆうやけ (yūyake)}]{\ja{ゆうやけ} (yūyake) — a sunset glow; a red evening sky}

\label{lesson:JA-C416-yuuyake}



\begin{quote}
Before the new one: say the Japanese for dawn, then the Japanese for sunrise.
\end{quote}

\subsection*{You'll want to know: \ja{ゆうやけ}}
\textbf{\ja{ゆうやけ}} — \emph{yūyake} — \textquotedblleft{}a sunset glow; a red evening sky\textquotedblright{}.

\begin{quote}
\textbf{\ja{うみの} \ja{ゆうやけは} \ja{きれいでした。}} — \emph{umi no yūyake wa kirei deshita} — \textquotedblleft{}The sunset over the sea was beautiful.\textquotedblright{}
\end{quote}

\begin{cousinweb}
\textbf{\ja{ゆうやけ}} (\emph{yūyake}) is a native Japanese word (wago), inherited rather than borrowed from Chinese, built from \emph{yū}, \textquotedblleft{}evening\textquotedblright{} and \emph{yake}, \textquotedblleft{}burning\textquotedblright{} (from \emph{yaku}); its kanji were attached later as kun'yomi (JMdict).
\end{cousinweb}

\begin{grammarlens}[title={What you've built}]
One more word for getting around.
\end{grammarlens}

\subsection*{Your turn}
\begin{itemize}
\raggedright
  \item \emph{Say it:} \emph{yūyake}
  \item \emph{Say it:} \emph{yūyake}, once more
  \item \emph{Say it:} \emph{umi no yūyake wa kirei deshita} — The sunset over the sea was beautiful.
  \item \emph{Recall:} say the Japanese for dawn, then the Japanese for sunrise, then say \emph{yūyake} again
\end{itemize}

\begin{tcolorbox}[breakable,colback=teal!4,colframe=teal!35!black,title={Before you move on}]
What does \ja{ゆうやけ} mean? (\textquotedblleft{}A sunset glow; a red evening sky\textquotedblright{}.) Say it once more.
\end{tcolorbox}
\section[\texorpdfstring{\ja{ひえこむ} (hiekomu)}{ひえこむ (hiekomu)}]{\ja{ひえこむ} (hiekomu) — to turn very cold}

\label{lesson:JA-C416-hiekomu}



\begin{quote}
Before the new one: say the Japanese for sunrise, then the Japanese for a sunset glow.
\end{quote}

\subsection*{You'll want to know: \ja{ひえこむ}}
\textbf{\ja{ひえこむ}} — \emph{hiekomu} — \textquotedblleft{}to turn very cold\textquotedblright{}.

\begin{quote}
\textbf{\ja{ふゆの} \ja{あさは} \ja{ひえこみます。}} — \emph{fuyu no asa wa hiekomimasu} — \textquotedblleft{}Winter mornings turn very cold.\textquotedblright{}
\end{quote}

Polite \textbf{\ja{ひえこみます}}, \emph{hiekomimasu}; a godan verb. Past (polite): \textbf{\ja{けさは} \ja{ひえこみました。}}, \emph{kesa wa hiekomimashita}.

\begin{cousinweb}
\textbf{\ja{ひえこむ}} (\emph{hiekomu}) is a native Japanese word (wago), inherited rather than borrowed from Chinese, built from \emph{hie}, \textquotedblleft{}growing cold\textquotedblright{} (from \emph{hieru}) and \emph{komu}, \textquotedblleft{}going deep\textquotedblright{}; its kanji were attached later as kun'yomi (JMdict).
\end{cousinweb}

\begin{grammarlens}[title={What you've built}]
One more word for getting around.
\end{grammarlens}

\subsection*{Your turn}
\begin{itemize}
\raggedright
  \item \emph{Say it:} \emph{hiekomu}
  \item \emph{Say it:} \emph{hiekomu}, once more
  \item \emph{Say it:} \emph{fuyu no asa wa hiekomimasu} — Winter mornings turn very cold.
  \item \emph{Recall:} say the Japanese for sunrise, then the Japanese for a sunset glow, then say \emph{hiekomu} again
\end{itemize}

\begin{tcolorbox}[breakable,colback=teal!4,colframe=teal!35!black,title={Before you move on}]
What does \ja{ひえこむ} mean? (\textquotedblleft{}To turn very cold\textquotedblright{}.) Say it once more.
\end{tcolorbox}
\section[\texorpdfstring{\ja{てらす} (terasu)}{てらす (terasu)}]{\ja{てらす} (terasu) — to shine on; to light up}

\label{lesson:JA-C416-terasu}



\begin{quote}
Before the new one: say the Japanese for a sunset glow, then the Japanese for to turn very cold.
\end{quote}

\subsection*{You'll want to know: \ja{てらす}}
\textbf{\ja{てらす}} — \emph{terasu} — \textquotedblleft{}to shine on; to light up\textquotedblright{}.

\begin{quote}
\textbf{\ja{つきが} \ja{みちを} \ja{てらしています。}} — \emph{tsuki ga michi o terashite imasu} — \textquotedblleft{}The moon is lighting up the road.\textquotedblright{}
\end{quote}

Polite \textbf{\ja{てらします}}, \emph{terashimasu}; a godan verb. Past (polite): \textbf{\ja{あさひが} \ja{やまを} \ja{てらしました。}}, \emph{asahi ga yama o terashimashita}.

\begin{cousinweb}
\textbf{\ja{てらす}} (\emph{terasu}) is a native Japanese word (wago), inherited rather than borrowed from Chinese; its kanji was attached later as a kun'yomi (JMdict).
\end{cousinweb}

\begin{grammarlens}[title={What you've built}]
One more word for getting around.
\end{grammarlens}

\subsection*{Your turn}
\begin{itemize}
\raggedright
  \item \emph{Say it:} \emph{terasu}
  \item \emph{Say it:} \emph{terasu}, once more
  \item \emph{Say it:} \emph{tsuki ga michi o terashite imasu} — The moon is lighting up the road.
  \item \emph{Recall:} say the Japanese for a sunset glow, then the Japanese for to turn very cold, then say \emph{terasu} again
\end{itemize}

\begin{tcolorbox}[breakable,colback=teal!4,colframe=teal!35!black,title={Before you move on}]
What does \ja{てらす} mean? (\textquotedblleft{}To shine on; to light up\textquotedblright{}.) Say it once more.
\end{tcolorbox}
